\section{Văn bản đầu tiên}
\label{sec:van-ban}

\subsection{Tạo dự án}
\begin{enumerate}
  \item Tại trang \nut{Projects}, bạn nhấp vào nút \nut{New project}.
  \item Chọn \nut{Blank project} để tạo dự án trống, hoặc chọn một mẫu có sẵn.
  \item Nhập tên dự án rồi bấm \nut{Create}, không gian làm việc sẽ mở ra ngay
    lập tức.
\end{enumerate}

\subsection{Khung của một tài liệu}
Mọi tài liệu \LaTeX{} đều gồm hai phần chính: \emph{phần đầu} (nằm trước lệnh
\lenh{begin\{document\}}) dùng để khai báo cấu hình, và \emph{phần thân} là nơi
chứa nội dung bạn viết.

\begin{ma}
\documentclass{article}
\usepackage[T5]{fontenc}
\usepackage[utf8]{inputenc}
\usepackage[vietnamese]{babel}

\title{Báo cáo của tôi}
\author{Nguyễn Văn A}

\begin{document}
\maketitle

\section{Mở đầu}
Đây là đoạn văn đầu tiên.

\end{document}
\end{ma}

\begin{chuy}
Ba dòng \lenh{usepackage} ở trên là bắt buộc để soạn thảo \textbf{tiếng Việt}.
Nếu thiếu chúng, trình biên dịch sẽ báo lỗi ``Unicode character'' tại các ký tự
có dấu (xem cách xử lý ở chương~\ref{sec:loi}).
\end{chuy}

\subsection{Biên dịch và xem PDF}
Để xem kết quả, bạn nhấn \phim{Ctrl}+\phim{S} (hoặc bấm nút ▶). Sau vài giây,
bản PDF sẽ xuất hiện ở khung bên phải. Mỗi khi chỉnh sửa nội dung, bạn chỉ cần
biên dịch lại để cập nhật bản xem trước.

\subsection{Chữ đậm, chữ nghiêng}
\begin{itemize}
  \item Dùng \lenh{textbf\{chữ đậm\}} (hoặc phím tắt \phim{Ctrl}+\phim{B}) để
    tạo \textbf{chữ đậm}.
  \item Dùng \lenh{emph\{nhấn mạnh\}} (hoặc phím tắt \phim{Ctrl}+\phim{I}) để
    in nghiêng \emph{nhấn mạnh}.
  \item Để ngắt sang đoạn mới, bạn cần chừa \textbf{một dòng trống} giữa hai
    đoạn, vì việc xuống dòng một lần trong mã nguồn sẽ không tạo ra đoạn mới.
\end{itemize}

\begin{meo}
Nếu không nhớ chính xác tên lệnh, bạn chỉ cần gõ dấu \texttt{\textbackslash}
kèm vài chữ cái đầu để SumiTeX tự động gợi ý. Chẳng hạn, khi gõ
\texttt{\textbackslash fig} rồi nhấn \phim{Enter}, hệ thống sẽ chèn sẵn một
khung hình mẫu và tự thêm gói \texttt{graphicx} nếu tài liệu chưa có.
\end{meo}

\begin{thuchanh}[Bài 1]
Mở tệp \texttt{bai-01.tex}, điền tên của bạn vào mục \lenh{author}, sau đó thêm
một mục \lenh{section\{Giới thiệu\}} chứa một đoạn văn ngắn có ít nhất một từ
in đậm rồi biên dịch để xem kết quả.
\end{thuchanh}